\chapter{Introduction}
\label{ch:intro}

Software security depends on finding vulnerabilities in source code before an
attacker does. Automated vulnerability detection tools aim to do this at scale,
flagging the lines of code where an injection, a traversal, or an unsafe
deserialization can occur. The accuracy of such a tool is only as trustworthy
as the benchmark used to measure it, and most current benchmarks inherit their
labels from records that were never written to mark code as vulnerable. This
chapter introduces the project domain of automated Python vulnerability
detection, sets out its global and societal relevance, identifies the gaps that
motivate the work, and states the objectives, scope, methodology, and
organization of the report.

%----------------------------------------------------------

\section{Overview of Project Topic}

The project builds and evaluates a benchmark for Python vulnerability
detection. A benchmark in this setting is a labeled corpus of code samples
paired with a scoring protocol, against which detectors are run so that their
results can be compared on common ground. The work covers two axes that prior
benchmarks treat in isolation: the correctness of the labels, and the
robustness of a detector when the input code is rewritten without changing what
it does.

\subsection{Automated Vulnerability Detection}

A vulnerability detector reads source code and predicts whether it contains a
security flaw, and which class of flaw. The classes follow the Common
Weakness Enumeration (CWE) catalog, for example SQL injection (CWE-89) or
cross-site scripting (CWE-79). Detectors fall into distinct families. Pattern
matching static application security testing (SAST) tools such as Bandit and
Semgrep scan a single file for known unsafe constructs. Whole-project dataflow
tools such as CodeQL and Snyk Code trace data from an input source to a
dangerous sink across an entire application. Large language models (LLMs)
reason over code in natural language, and learned classifiers such as
GraphCodeBERT are fine-tuned on labeled vulnerability data. Each family makes a
different trade-off between coverage, precision, and the context it requires.

\subsection{Global Scenario}

Most research in this area is built on C and C++. The standard datasets,
BigVul, CVEfixes, CrossVul, DiverseVul, Devign, and ReVeal, are dominated by C
and C++ code and contribute few Python samples per CWE class. Python is the
dominant language for security tooling, machine learning pipelines, and modern
web applications, yet it has no curated, audit-validated detection benchmark
comparable to the CASTLE benchmark that exists for C and C++. The result is a
mismatch between where security-sensitive Python code is written and where
detection methods are measured.

\subsection{Societal Relevance}

Web applications and data pipelines written in Python process personal data,
financial records, and critical infrastructure traffic. A single missed
injection or deserialization flaw can expose an entire user base. Detection
tools shift the cost of finding these flaws from the attacker to the defender,
but only if their reported accuracy reflects real detection ability rather than
agreement with mislabeled data. A benchmark with validated labels and a
robustness axis gives practitioners a more honest basis for choosing a tool,
which raises the security floor of the software that society depends on.

\subsection{Problem Area}

Two problems define the gap this project addresses. First, the labels in
existing benchmarks are noisy. They are inherited from CVE fix commits, and a
fix commit touches several files for refactoring, tests, or cleanup that have
nothing to do with the vulnerability. Each co-changed file inherits
the CWE label of the advisory, and manual audits report between 25\% and 50\%
label noise as a result \cite{croft,cleanvul}. A benchmark built on raw
inheritance therefore measures a detector's agreement with label noise as much
as its detection ability \cite{primevul}. Second, accuracy is reported on
unmodified code only. A reported score is an upper bound that may not survive a
trivial, behavior-preserving rewrite of the input \cite{secllmholmes}, so it
says little about how a detector behaves on code it has not seen in that exact
form.

%----------------------------------------------------------

\section{Specific Details of the Project Topic}

The benchmark covers seven \emph{sink-shaped} CWE classes drawn from the MITRE
Top-25 \cite{mitre_top25}: SQL injection (CWE-89), cross-site scripting
(CWE-79), path traversal (CWE-22), OS command injection (CWE-78), code
injection (CWE-94), server-side request forgery (CWE-918), and insecure
deserialization (CWE-502). A sink-shaped class has a well-defined unsafe
function, the sink, whose misuse is the vulnerability and whose fix modifies
the sink call. Structural classes that encode the absence of a check, such as
missing authorization, require application-level context and are excluded on
methodological grounds.

The corpus holds 869 samples: 440 real-world positives validated by audit and
429 synthesized safe hard negatives. Labels pass through a three-layer
validation methodology, and the corpus is paired with seven semantics-preserving
mutators that probe detector shortcuts. Seven detectors across four tiers are
evaluated on both clean and mutated code.

\section{Literature Review}

The literature divides into four threads that this project unifies. Large
vulnerability datasets establish the supply of labeled code but are C and C++
centric \cite{bigvul,diversevul}. Studies of label noise, by Croft et al., the
ReVeal authors, CleanVul, and PrimeVul, quantify how inherited CVE labels
mislead detector evaluation \cite{croft,reveal,cleanvul,primevul}. Comparative
studies such as CASTLE pit static analyzers against LLMs but rely on
CVE-inherited labels and add no robustness axis \cite{castle}. Robustness
studies, including SecLLMHolmes, show that semantics-preserving edits can flip a
detector's verdict \cite{secllmholmes}. Chapter~\ref{ch:theory} reviews these
threads in detail.

\section{Problem Statement}

To construct an audit-validated benchmark for Python vulnerability detection
across sink-shaped CWE classes, with a label-validation methodology that reduces
inherited label noise, and to measure detector accuracy and adversarial
robustness jointly so that the evaluation distinguishes detectors that reason
about code from detectors that pattern-match its surface.

\section{Objectives}

\begin{itemize}
  \item Assemble a Python vulnerability corpus across seven sink-shaped CWE
    classes from advisory and CVE-fix sources, paired with constructed safe
    hard negatives.
  \item Design and apply a three-layer label-validation methodology that
    lowers the audited false-positive rate of the labels.
  \item Build a suite of semantics-preserving mutators, each targeting one
    detector shortcut, to probe robustness.
  \item Evaluate seven detectors across four tiers on clean and mutated code
    under a set of bounded, complementary metrics.
  \item Release all artifacts, audit reports, and rejection manifests to
    support replication.
\end{itemize}

\section{Scope}

The benchmark is Python-specific and limited to sink-shaped CWE classes. It does
not cover structural CWEs, multi-language code, or whole-application analysis,
and the scores are not a comprehensive ranking of Python security tools. The
validation methodology is language-agnostic in design, but its sink pattern set
would need re-deriving for another language. The trained detector is presented
as a small-data existence proof rather than a state-of-the-art claim.

\section{Methodology}

The work proceeds in four stages. First, positives are scraped from five
advisory and CVE-fix sources and passed through a sink-presence filter, an
independent adversarial audit, and a diff-based filter that anchors each
positive on a modified sink line. Second, safe hard negatives are constructed as
remediated twins of real vulnerable functions. Third, seven mutators are
implemented over the Python abstract syntax tree (AST), each deterministic and
semantics-preserving. Fourth, seven detectors are run on the held-out test set,
clean and mutated, and scored under five bounded metrics. Each stage is
reproducible from the released artifacts.

\section{Organization of the Report}

Chapter~\ref{ch:theory} presents the theory and concepts behind vulnerability
detection, label noise, and robustness, together with related work.
Chapter~\ref{ch:srs} specifies the functional and non-functional requirements of
the benchmark system. Chapter~\ref{ch:hld} gives the high-level design,
including the system architecture and data flow. Chapter~\ref{ch:detailed}
details the design of each module. Chapter~\ref{ch:impl} describes the
implementation, including the dataset pipeline, mutators, and trained detector.
Chapter~\ref{ch:testing} covers the testing of the pipeline and harness.
Chapter~\ref{ch:results} reports the experimental results across the seven
detectors. Chapter~\ref{ch:conclusion} concludes with limitations and future
enhancements.
